\begin{abstractpage}

\begin{abstract}{romanian}
Detecția de anomalii în seriile temporale joacă un rol critic în asigurarea securității pentru diferite sisteme informatice, printre care sisteme industriale, economice și de sănătate. În ultimii ani, metodele de învățare profundă cu rețele neuronale pentru a prezice și reconstrui serii temporale au devenit din ce în ce mai populare datorită abilității lor de a modela structuri temporale complexe. În timp ce rețelele generative adversariale (GAN) au avut parte de un mare succes în vederea artificială, transferarea performanțelor către domeniul detecției de anomalii în seriile temporale este o zonă de cercetare activă. În această disertație, noi evaluăm performanța unui GAN standard și a unui GAN ciclic, în comparație cu performanța obținută de un LSTM folosit pentru prezicere, cât și pentru reconstrucție. Evaluarea rețelelor a fost făcută folosind trei seturi de date: Numenta Anomaly Benchmark (NAB), Yahoo S5 și setul de date cu telemetrii de la NASA. În urma rezultatelor, am constatat că integrarea unei penalizări ciclice îmbunătățește performanța modelelor GAN, confirmând faptul că arhitecturile bazate pe GAN reprezintă o direcție promițătoare pentru cercetarea detecției de anomalii în serii temporale.

\end{abstract}

\begin{abstract}{english}
Time-series anomaly detection plays a critical role in ensuring the security of multiple systems, including industrial, economic, and healthcare environments. Concurrently, deep learning methods for time-series forecasting and reconstruction have grown significantly in popularity due to their ability to model complex temporal patterns. While Generative Adversarial Networks (GANs) have achieved massive success in computer vision, transferring their performance to the time-series domain remains an active area of exploration. In this thesis, we evaluate the performance of a standard GAN and a cycle GAN against LSTM-based prediction and reconstruction. Evaluation of these approaches was made across three benchmark datasets: Numenta Anomaly Benchmark (NAB), the Yahoo S5, and the NASA telemetry datasets. Our results show that incorporating a cycle loss improves GAN performance, confirming that GAN-based architectures represent a promising direction for future time-series anomaly detection research.
\end{abstract}

\end{abstractpage}